\newpage
\section{Main theorem}

\subsection{Notations \& Énoncé}

\begin{lemme}
Il existe $\underline{C}^\ensddfv$, $\overline{C}^\ensddfv$, $\underline{C}^\ensdiamant$ et $\overline{C}^\ensdiamant$ indépendantes de $h$ et dépendantes de $\reg \ensddfv$ (refaire la démo de \cite{dziuk_finite_2013} pour bien comprendre de quoi dépendent ces constantes) telles que
\[
\forall \primal{K} \in \ensprimal,
\quad
2 \, \underline{C}^\ensddfv \, \poids{\primal{K}} \leqslant \mesure(\primal{K}) \leqslant 2 \, \overline{C}^\ensddfv \, \poids{\primal{K}},
\]
\[
\forall \mailleduale{K} \in \ensdual,
\quad
2 \, \underline{C}^\ensddfv \, \poids{\mailleduale{K}} \leqslant \mesure(\mailleduale{K}) \leqslant 2 \, \overline{C}^\ensddfv \, \poids{\mailleduale{K}},
\]
\[
\forall \diamant{D} \in \ensdiamant,
\quad
\underline{C}^\ensdiamant \, \poids{\diamant{D}} \leqslant \mesure(\diamant{D}) \leqslant \overline{C}^\ensdiamant \, \poids{\diamant{D}}.
\]
\end{lemme}

\begin{prop}
\[
\normeddfvD{} \sim \norme{}_{\L^2(\variete{M})}   
\]
\end{prop}

\begin{prop}
\begin{equation}
\underline{C}^\ensddfv \leqslant \norme{\projddfvvar}_{\L^2} \leqslant \overline{C}^\ensddfv,
\quad
\underline{C}^\ensdiamant \leqslant \norme{\projdiamantvar}_{\L^2} \leqslant \overline{C}^\ensdiamant.
\end{equation}
\end{prop}

\begin{demo}
\[
\norme{\projddfvvar}_{\L^2} = \sup_{u^\ensddfv \in \R^\ensddfv} \frac{\norme{\projddfvvar u^\ensddfv}_{\L^2(\variete{M})}}{\normeddfvT{u^\ensddfv}}
\]
\end{demo}

Nous sommes maintenant en mesure d'énoncer le résultat principal de ce chapitre, à savoir une inégalité de Poincaré \ddfv sur une surface fermée sans bord de $\R^3$.
\begin{theo}[Main]
    Soit $\ensddfv$ un maillage \ddfv de $\variete{M}$. Il existe une constante $C$, ne dépendant que de la régularité du maillage $\reg(\ensddfv)$, telle que pour tout $u^\ensddfv \in \R^\ensddfv$, on a
    \begin{equation} \label{eq:main_ddfv}
    \normeddfvT{u^\ensddfv} \leqslant C \normeddfvD{\grad^\ensdiamant u^\ensddfv}.
    \end{equation}
    De manière équivalente ?, on a
    \[
    \norme{\projddfvvar u^\ensddfv}_{\L^2(\variete{M})} \leqslant C \, \frac{\overline{C}^\ensdiamant}{\underline{C}^\ensddfv} \norme{\projdiamantvar \grad^\ensdiamant u^\ensddfv}_{\L^2(\variete{M})}.
    \]
\end{theo}

\subsection{Une inégalité de graphe}

\begin{defi} \label{defi:graphe_ddfv}
Soit $\ensddfv= (\ensprimal, \ensdual)$ un maillage \ddfv. On lui associe deux graphes:
\begin{itemize}
    \item le \defini{graphe du maillage primal} $\graphe^\ensprimal = (\enssomdual, \ensaretedual)$ où $\enssomdual$ est l'ensemble des sommets de $\ensdual$ et $\ensaretedual$ est l'ensemble des arêtes de $\ensdual$;
    \item le \defini{graphe du maillage dual} $\graphe^\ensdual = (\enssomprim, \ensareteprim)$ où $\enssomprim$ est l'ensemble des sommets de $\ensprimal$ et $\ensareteprim$ est l'ensemble des arêtes de $\ensprimal$.
\end{itemize}
\end{defi}

\begin{envide}[Notations]
On note $\vertices(\graphe)$ l'ensemble des sommets d'un graphe $\graphe$ et $\edges(\graphe)$ l'ensemble de ses arêtes. On pourra également utiliser la notation $\vertices$ et $\edges$ lorsqu'il n'y a pas d'ambiguïté sur le graphe. 

Soit $u^\ensddfv = \big( u^\ensprimal, u^\ensdual \big) \in \R^\ensddfv$. On considère la fonction $u^\ensprimal$ comme une fonction sur la graphe $\graphe^\ensprimal$ dans le sens où pour $x_{\primal{K}} \in \enssomdual$, on note $u_{\primal{K}} = u^\ensprimal(x_{\primal{K}})$ et on considère la fonction $u^\ensdual$ comme une fonction sur la graphe $\graphe^\ensdual$ dans le sens où pour $x_{\mailleduale{K}} \in \enssomprim$, on note $u_{\mailleduale{K}} = u^\ensdual(x_{\mailleduale{K}})$.
\end{envide}

\begin{figure}[H]
    \centering
    \tikzset{
    ddfv/.style={
        scale=1.5,
        every path/.style={thick,line join=round, cap=round},
        ptsq/.style 2 args={
          rectangle,
          fill=##1,
          inner sep=2pt,
          label={[text=##1]##2}
        },
        ptsqbord/.style 2 args={
          rectangle,
          draw=##1,
          inner sep=2pt,
          label={[text=##1]##2}
        },
        ptcl/.style 2 args={
          circle,
          fill=##1,
          inner sep=1.5pt,
          label={[text=##1]##2}
        },
        ptclbord/.style 2 args={
          circle,
          draw=##1,
          inner sep=1.5pt,
          label={[text=##1]##2}
        },
        mailleprimale/.style={
            myblue,
            fill=mylightblue
        },
        mailledualeint/.style={
            myred,
            fill=mylightred
        },
        mailledualebord/.style={
            myred,
            fill=mylightred!50!white
        },
        maillediamantint/.style={
            mygreen,
            fill=mylightgreen
        },
        maillediamantbord/.style={
            mygreen,
            fill=mylightgreen!50!white
        }
    }
}

\begin{tikzpicture}[ddfv]

% \foreach \i in {1,3,4,5,8,12,13,14,16,17,18,19}


\coordinate (M1)  at (-0.2,-0.3); 
\coordinate (M1b) at (-0.4,-0.2); 
\coordinate (M2)  at (1.1,0.2);
\coordinate (M3)  at (2.0,-0.3); 
\coordinate (M3b) at (2.2,-0.3); 
\coordinate (M4)  at (3.2,-0.3); 
\coordinate (M5)  at (-0.4,0.9); 
\coordinate (M5b) at (-0.4,1.1); 
\coordinate (M6)  at (0.8,1.2);
\coordinate (M7)  at (1.9,0.9);
\coordinate (M8)  at (3.2,1.1);
\coordinate (M9)  at (0.1,2.1);
\coordinate (M10) at (1.0,2.0);
\coordinate (M11) at (2.2,2.3);
\coordinate (M12) at (3.2,1.8);
\coordinate (M12b) at (3.2,2.2);
\coordinate (M13)  at (-0.4,3.2); 
\coordinate (M14)  at (0.9,3.2);
\coordinate (M14b) at (0.7,3.2);
\coordinate (M14t) at (1.1,3.2);
\coordinate (M15)  at (2.0,2.8);
\coordinate (M16)  at (3.1,3.2); 
\coordinate (M16b) at (3.2,3.1); 
\coordinate (M17)  at (1,-0.3);
\coordinate (M18)  at (2,3.2);
\coordinate (M19)  at (-0.4,1.8);

\coordinate (Me1)  at (0.2,0.7);
\coordinate (Me2)  at (0.7,0.3);
\coordinate (Me3)  at (1,-0.1);
\coordinate (Me4)  at (2.7,0.9);
\coordinate (Me5)  at (0.2,1.4);
\coordinate (Me6)  at (1.3,0.8);
\coordinate (Me7)  at (2.5,1.8);
\coordinate (Me8)  at (0.7,2.4);
\coordinate (Me9)  at (1.4,2.5);
\coordinate (Me10) at (2.8,2.4);
\coordinate (Me11) at (1.2,1.4);
\coordinate (Me12) at (1.8,1.6);
\coordinate (Me13) at (2.3,0.2);
\coordinate (Me14) at (0.6,1.8);
\coordinate (Me15) at (1.7,3.0);
\coordinate (Me16) at (2.5,2.8);
\coordinate (Me17) at (-0.1,2.3);
\coordinate (Me18) at (2.3,3.0);
\coordinate (Me19) at (-0.2,1.7);
\coordinate (Me20) at (0.3,2.6);
\coordinate (Me21) at (1.8,0.4);

\coordinate (Mb1) at ($(M1b)!0.5!(M5)$);
\coordinate (Mb2) at ($(M19)!0.5!(M5)$);
\coordinate (Mb3) at ($(M19)!0.5!(M13)$);
\coordinate (Mb4) at ($(M13)!0.5!(M14)$);
\coordinate (Mb5) at ($(M14)!0.5!(M18)$);
\coordinate (Mb6) at ($(M16)!0.5!(M18)$);
\coordinate (Mb7) at ($(M16b)!0.5!(M12)$);
\coordinate (Mb8) at ($(M12)!0.5!(M8)$);
\coordinate (Mb9) at ($(M4)!0.5!(M8)$);
\coordinate (Mb10) at ($(M4)!0.5!(M3)$);
\coordinate (Mb11) at ($(M17)!0.5!(M3)$);
\coordinate (Mb12) at ($(M17)!0.5!(M1)$);

% \draw[mailleprimale] (M1) -- (M5) -- (M19) -- (M13) -- (M14) -- (M18) -- (M16) -- (M12) -- (M8) -- (M4) -- (M3) -- (M17) -- cycle;

% \draw (-0.4,-0.3) rectangle (M16);

\foreach \i in {2,6,7,9,10,11,15} {
   % \node[ptsq={myblue}{above:{\i}}] at (M\i) {};
   \node[ptsq={myblue}{}] at (M\i) {};
}

\foreach \i in {1,3,4,5,8,12,13,14,16,17,18,19} {
   % \node[ptsqbord={myblue}{above:{\i}}] at (M\i) {};
   % \node[ptsq={myblue}{}] at (M\i) {};
}

\foreach \i in {1,...,21} {
   % \node[ptcl={myred}{above:{\i}}] at (Me\i) {};
   \node[ptcl={myred}{}] at (Me\i) {};
}

\draw[myblue] (M16) -- (M15) -- (M14t);
\draw[myblue] (M1b) -- (M6) -- (M9) -- (M14b);
\draw[myblue] (M14) -- (M10) -- (M11) -- (M15) -- (M18);
\draw[myblue] (M11) -- (M12b);
\draw[myblue] (M12) -- (M7) -- (M6)  -- (M10) -- (M9) -- (M19);
\draw[myblue] (M2) -- (M6) -- (M5); % -- (M19) -- (M13);
\draw[myblue] (M11) -- (M16b);
\draw[myblue] (M10) -- (M7) -- (M2);
\draw[myblue] (M4) -- (M7);
\draw[myblue] (M5b) -- (M9) -- (M13);
\draw[myblue] (M1) -- (M2) -- (M3);
\draw[myblue] (M11) -- (M7) -- (M3b);


% \foreach \i in {1,...,12} {
%     % \node[ptcl={myred}{above:{\i}}] at (Mb\i) {};
%     \node[ptclbord={myred}{}] at (Mb\i) {};
% }


\draw[myred] (Me1) -- (Me2) -- (Me6) -- (Me11) -- (Me14);
\draw[myred] (Me2) -- (Me3) -- (Me21);
\draw[myred] (Me21) -- (Me13) -- (Me4) -- (Me7) -- (Me12) -- (Me11) -- (Me6) -- cycle;
\draw[myred] (Me14) -- (Me11) -- (Me12) -- (Me9) -- (Me8);
\draw[myred] (Me10) -- (Me7) -- (Me12) -- (Me9);
\draw[myred] (Me5) -- (Me14) -- (Me8) -- (Me20) -- (Me17);
\draw[myred] (Me9) -- (Me16) -- (Me18);

\draw[myred] (Mb11) -- (Me3) -- (Mb12);
\draw[myred] (Mb12) -- (Me3) -- (Me2) -- (Me1) -- (Mb1);
\draw[myred] (Mb1) -- (Me1) -- (Me5) -- (Me19) -- (Mb2);
\draw[myred] (Mb2) -- (Me19) -- (Me17) -- (Mb3);
\draw[myred] (Mb3) -- (Me17) -- (Me20) -- (Mb4);
\draw[myred] (Mb4) -- (Me20) -- (Me8) -- (Me9) -- (Me15) -- (Mb5);
\draw[myred] (Mb5) -- (Me15) -- (Me18) -- (Mb6);
\draw[myred] (Mb6) -- (Me18) -- (Me16) -- (Me10) -- (Mb7);
\draw[myred] (Mb7) -- (Me10) -- (Me7) -- (Me4) -- (Mb8);
\draw[myred] (Mb8) -- (Me4) -- (Mb9);
\draw[myred] (Mb9) -- (Me4) -- (Me13) -- (Mb10);
\draw[myred] (Mb10) -- (Me13) -- (Me21) -- (Me3) -- (Mb11);

\node[myred, below=5pt] at (Me12) {\contour{white}{$x_{\primal{K}}$}};
\node[myred, above=5pt] at (Me9) {\contour{white}{$x_{\primal{L}}$}};

\node[myred] at ($(Me12)!0.5!(Me9)$) {\contour{white}{$\dual{\sigma}$}};

\node[myblue, shift={(2pt,15pt)}] at (Me12) {\contour{white}{$\primal{K}$}};
\node[myblue, shift={(17pt,-5pt)}] at (Me9) {\contour{white}{$\primal{L}$}};


\end{tikzpicture}
    \caption{Maillage primal et graphe associé}
\end{figure}

\noindent Pour tout $u^\ensddfv \in \R^\ensddfv$, on pose
\[
\mathcal{E}(u^\ensddfv) \defeq 
\sum_{\primal{K} \sim \primal{L}} ( u_{\primal{K}} - u_{\primal{L}} )^2
+
\sum_{\mailleduale{K} \sim \mailleduale{L}} ( u_{\mailleduale{K}} - u_{\mailleduale{L}} )^2.
\]
Pour obtenir \Cref{eq:main_ddfv}, l'objectif va être d'obtenir d'une part $\mathcal{E}(u^\ensddfv) \gtrsim \normeddfvT{u^\ensddfv}^2$ et d'autre part $\mathcal{E}(u^\ensddfv) \lesssim \normeddfvD{\grad^\ensdiamant u^\ensddfv}^2$. 

\begin{lemme}[Une inégalité de graphe] \label{lemme:ineg_de_graphe}
Pour tout $u^\ensddfv \in \R^\ensddfv$,
\[
\mathcal{E}(u^\ensddfv) \geqslant
\frac{h_{\graphe^\ensddfv}{}^2}{2} \mathopen{}\left(
    \sum_{\primal{K} \in \ensprimal} \big( u_{\primal{K}} - c^\ensprimal \big)^2 \degr(x_{\primal{K}})
    +
    \sum_{\mailleduale{K} \in \ensdual} \big( u_{\mailleduale{K}} - c^\ensdual \big)^2 \degr(x_{\mailleduale{K}})
\right),
\]
où
\[
c^\ensprimal = \frac{
    \sum\limits_{\primal{K} \in \ensprimal} u_{\primal{K}} \degr(x_{\primal{K}})
}{
    \sum\limits_{\primal{K} \in \ensprimal} \degr(x_{\primal{K}})
},
\quad
c^\ensdual = \frac{
    \sum\limits_{\mailleduale{K} \in \ensdual} u_{\mailleduale{K}} \degr(x_{\mailleduale{K}})
}{
    \sum\limits_{\mailleduale{K} \in \ensdual} \degr(x_{\mailleduale{K}})
},
\quad
\text{et}
\quad
h_{\graphe^\ensddfv} = \min\mathopen{}\big\{ h_{\graphe^\ensprimal}, h_{\graphe^\ensdual}\big\}.
\]
La constante $h_{\graphe^\ensprimal}$ (resp. $h_{\graphe^\ensdual}$) est la constante de Cheeger du graphe $\graphe^\ensprimal$ (resp. $\graphe^\ensdual$) définie par \Cref{eq:cheeger_constant}.
\end{lemme}

\begin{demo}
Soit $u^\ensddfv \in \R^\ensddfv$. On considère cette fonction comme une fonction sur les graphes $\graphe^\ensprimal$ et $\graphe^\ensdual$ comme expliqué dans la \Cref{defi:graphe_ddfv}. Il suffit alors d'appliquer le \Cref{theo:cheeger_inequality}, établissant l'inégalité de Cheeger, d'une part au graphe $\graphe^\ensprimal$ : 
\begin{equation} \label{eq:ineg_de_graphe_primal}
\sum_{\primal{K} \sim \primal{L}} ( u_{\primal{K}} - u_{\primal{L}} )^2
\geqslant 
\frac{h_{\graphe^\ensprimal}{}^2}{2}
\sum_{\primal{K} \in \ensprimal} \big( u_{\primal{K}} - c^\ensprimal \big)^2 \degr(x_{\primal{K}}),
\end{equation}
et d'autre part au graphe $\graphe^\ensdual$:
\begin{equation} \label{eq:ineg_de_graphe_dual}
\sum_{\mailleduale{K} \sim \mailleduale{L}} ( u_{\mailleduale{K}} - u_{\mailleduale{L}} )^2
\geqslant 
\frac{h_{\graphe^\ensdual}{}^2}{2}
\sum_{\mailleduale{K} \in \ensdual} \big( u_{\mailleduale{K}} - c^\ensdual \big)^2 \degr(x_{\mailleduale{K}}),
\end{equation}
où on a utilisé l'expression \Cref{eq:spectral_gap} pour les \emph{spectal gaps}. La somme des inégalités \ref{eq:ineg_de_graphe_primal} et \ref{eq:ineg_de_graphe_dual} donne le résultat. 
\end{demo}

\begin{lemme} \label{lemme:minoration_degre}
Soient $\primal{K} \in \ensprimal$ et $\mailleduale{K} \in \ensdual$. Alors
\[
\degr(x_{\primal{K}}) \geqslant \frac{2 \mesure(\primalplat{K})}{\reg(\ensddfv)^2 \size(\ensddfv)^2},
\quad
\degr(x_{\mailleduale{K}}) \geqslant \frac{2 \mesure(\mailledualeplate{K})}{\reg(\ensddfv)^2 \size(\ensddfv)^2}.
\]
\end{lemme}

\begin{demo}
On raisonne sur une maille primale plate $\primalplat{K}$. Posons $n_{\primal{K}} \defeq \degr(x_{\primal{K}})$. Chaque arête de $\primalplat{K}$ correspond à une maille primale voisine donc $n_{\primal{K}}$ est le nombre d'arêtes de $\primalplat{K}$. Puisque la maille est \textbf{étoilée} par rapport à $x_{\primal{K}}$, on peut la découper en $n_{\primal{K}}$ triangles d'intérieurs disjoints : 
\[
\primalplat{K}
=
\bigcup_{i = 1}^{n_{\primal{K}}}
\Delta(x_{\primal{K}}, s_i, s_{i+1}),
\]
où $(s_i)_{1 \leqslant i \leqslant n_{\primal{K}}}$ sont les sommets de la maille, ordonnés cycliquement avec la convention $s_{n_{\primal{K}} + 1} = s_1$. Si on note $\angle{i}$ l'angle $\widehat{s_i \, x_{\primal{K}} \, s_{i+1}}$, alors
\[
\mesure(\primalplat{K})
= 
\frac{1}{2} \sum_{i=1}^{n_{\primal{K}}}
\underbrace{\norme{s_i - x_{\primal{K}}}}_{\leqslant \diam{\primalplat{K}}} \underbrace{\norme{s_{i+1} - x_{\primal{K}}}}_{\leqslant \diam{\primalplat{K}}} \abs{\sin \angle{i}}
\leqslant
\frac{1}{2} \, n_{\primal{K}} \diam{\primalplat{K}}{}^2.
\] 
On obtient alors 
\begin{equation} \label{eq:estim}
\degr(x_{\primal{K}}) \geqslant \frac{2 \mesure(\primalplat{K})}{\diam{\primalplat{K}}{}^2}
\end{equation}
et le même raisonnement s'applique aux mailles duales plates : 
\[
\degr(x_{\mailleduale{K}}) \geqslant \frac{2 \mesure(\mailledualeplate{K})}{\diam{\mailledualeplate{K}}{}^2}.
\]
Dans la régularité $\reg(\ensddfv)$, le terme
\[
\max_{\primal{K} \in \ensprimal} \max_{\diamant{D} \in \ensdiamant_{\primal{K}}} \frac{\diam{\primalplat{K}}}{\diam{\diamant{D}}}
\]
permet d'avoir l'estimation
\begin{equation} \label{eq:estim_bis}
\diam{\primalplat{K}} \leqslant \reg(\ensddfv) \diam{\diamant{D}} {\color{red} \diam{\diamantplat{D}}} \leqslant \reg(\ensddfv) \size(\ensddfv)
\end{equation}
pour tout diamant incident à $\primalplat{K}$. La même propriété vaut pour les mailles duales. La combinaison de \Cref{eq:estim,eq:estim_bis} donne le résultat.
\end{demo}

\begin{remarque}
Avec les poids $\poids{\primal{K}} = \displaystyle\frac{\mesure(\primalplat{K})}{2}$ et $\poids{\mailleduale{K}} = \displaystyle\frac{\mesure(\mailledualeplate{K})}{2}$, les minorations deviennent
\[
\degr(x_{\primal{K}}) \geqslant \frac{4}{\reg(\ensddfv)^2 \size(\ensddfv)^2} \, \poids{\primal{K}},
\quad
\degr(x_{\mailleduale{K}}) \geqslant \frac{4}{\reg(\ensddfv)^2 \size(\ensddfv)^2} \, \poids{\mailleduale{K}}.
\]
\end{remarque}
Les \Cref{lemme:ineg_de_graphe,lemme:minoration_degre} vont nous permettre de montrer le résultat suivant. 

\begin{prop} \label{prop:minoration_energie}
Pour tout $u^\ensddfv \in \Rddfvmn$,
\[
\mathcal{E}(u^\ensddfv) \geqslant \frac{2 \, h_{\graphe^\ensddfv}{}^2}{\reg(\ensddfv)^2 \size(\ensddfv)^2} \normeddfvT{u^\ensddfv}^2.
\]
L'ensemble $\Rddfvmn$ défini par \Cref{eq:moyenne_nulle} est l'ensemble des fonctions scalaires \ddfv à moyenne nulle. 
\end{prop}

\begin{demo}
Soit $u^\ensddfv \in \Rddfvmn$. Par définition, la moyenne géométrique $\sum\limits_{\primal{K} \in \ensprimal} \poids{\primal{K}} \, u_{\primal{K}}$ est nulle. De plus, elle minimise la fonction $c \mapsto \sum\limits_{\primal{K} \in \ensprimal} \poids{\primal{K}} \, (u_{\primal{K}} - c)^2$. Par conséquent,
\[
\sum_{\primal{K} \in \ensprimal} \poids{\primal{K}} \big( u_{\primal{K}} - c^\ensprimal)^ 2
\geqslant
\sum_{\primal{K} \in \ensprimal} \poids{\primal{K}} \, u_{\primal{K}}{}^2,
\]
et en utilisant les estimations du \Cref{lemme:minoration_degre} on trouve
\begin{equation} \label{eq:ineg1}
\sum_{\primal{K} \in \ensprimal} \big(u_{\primal{K}} - c^\ensprimal \big)^2 \degr(x_{\primal{K}})
\geqslant
\frac{4}{\reg(\ensddfv)^2 \size(\ensddfv)^2}
\sum_{\primal{K} \in \ensprimal} \poids{\primal{K}} \, u_{\primal{K}}{}^2.
\end{equation}
Le même raisonnement sur le maillage dual donne
\begin{equation} \label{eq:ineg2}
\sum_{\mailleduale{K} \in \ensdual} \big(u_{\mailleduale{K}} - c^\ensdual \big)^2 \degr(x_{\mailleduale{K}})
\geqslant
\frac{4}{\reg(\ensddfv)^2 \size(\ensddfv)^2}
\sum_{\mailleduale{K} \in \ensdual} \poids{\mailleduale{K}} \, u_{\mailleduale{K}}{}^2.
\end{equation}
On somme les inégalités \ref{eq:ineg1} et \ref{eq:ineg2} puis le \Cref{lemme:ineg_de_graphe} et la définition de la norme $\normeddfvT{}$ donnent le résultat. 
\end{demo}

\noindent Intéressons-nous maintenant à la majoration de $\mathcal{E}$.

\begin{prop} \label{prop:majoration_energie}
Pour tout $u^\ensddfv \in \R^\ensddfv$,
\[
\mathcal{E}(u^\ensddfv)
\leqslant
4 \reg(\ensddfv)^2 \normeddfvD{\grad^\ensdiamant u^\ensddfv}^2.
\]
\end{prop}

\begin{demo}
Soient deux mailles primales voisines $\primal{K} \sim \primal{L}$ définissant le diamant $\diamant{D} = \diamant{D}_{\sigma', \dual{\sigma}}$. Par définition du gradient \ref{defi:discrete_gradient} sur $\diamant{D}$, on a
\begin{equation*}
    (u_{\primal{L}} - u_{\primal{K}})^2 
    = \mesure(\dual{\sigma})^2 \abs{\prodscal{\grad^{\diamant{D}} u^\ensddfv}{\vect{\tau}_{\primal{K},\primal{L}}}}^2
    \leqslant \mesure(\dual{\sigma})^2 \abs{\grad^{\diamant{D}} u^\ensddfv}^2.
\end{equation*}
De même, pour deux mailles duales voisines $\mailleduale{K} \sim \mailleduale{L}$ définissant le diamant $\diamant{D} = \diamant{D}_{\sigma,\dual{\sigma}'}$, on a
\begin{equation*}
    (u_{\mailleduale{L}} - u_{\mailleduale{K}})^2 
    = \mesure(\sigma)^2 \abs{\prodscal{\grad^{\diamant{D}} u^\ensddfv}{\vect{\tau}_{\mailleduale{K},\mailleduale{L}}}}^2
    \leqslant \mesure(\sigma)^2 \abs{\grad^{\diamant{D}} u^\ensddfv}^2.
\end{equation*}
Les ensembles $\ensareteprim$ et $\ensaretedual$ sont en bijection avec l'ensemble $\ensdiamant$. Ainsi
\begin{equation} \label{eq:inter}
\mathcal{E}(u^\ensddfv)
\leqslant
\sum_{\diamant{D} \in \ensdiamant} \big[\mesure(\sigma)^2 + \mesure(\dual{\sigma})^2\big] \abs{\grad^{\diamant{D}} u^\ensddfv}^2.
\end{equation}
On souhaite avoir $\mesure(\sigma)^2 + \mesure(\dual{\sigma})^2 \leqslant C \, \poids{\diamant{D}}$ avec $C > 0$ indépendante de $\size(\ensddfv)$. Comme 
\begin{equation} \label{eq:rapport}
\frac{
    \mesure(\sigma)^2 + \mesure(\dual{\sigma})^2 
}{
    \poids{\diamant{D}}
}
=
\frac{2}{\sin \angle{\diamant{D}}}
\left(
    \frac{\mesure(\sigma)}{\mesure(\dual{\sigma})}
    +
    \frac{\mesure(\dual{\sigma})}{\mesure(\sigma)}
\right),
\end{equation}
on met dans la régularité
\[
\max_{\diamant{D}_{\sigma, \dual{\sigma}} \in \ensdiamant}\mathopen{}\left\{
    \frac{\mesure(\sigma)}{\mesure(\dual{\sigma})},
    \frac{\mesure(\dual{\sigma})}{\mesure(\sigma)}
\right\}
\quad
\text{et}
\quad
\frac{1}{\sin \angle{\ensddfv}}.
\]
Autrement dit, il faut imposer que l'angle entre deux diagonales ne tende ni vers $0$ ni vers $\pi$ et que le rapport des longueurs des deux diagonales reste uniformément borné. 

\begin{remarque}
Ceci est contrôlé par
\[
    \max_{\diamant{D}\in\ensdiamant}
    \max_{\quart{Q}\in\ensquart_{\diamant{D}}}
    \frac{\diam{\diamant{D}}}
    {\min\limits_{\eta\in\partial\quart{Q}}\mesure(\eta)}.
\]
\end{remarque}
Ainsi, d'après \Cref{eq:rapport}, 
\begin{align*}
\mesure(\sigma)^2 + \mesure(\dual{\sigma})^2 
\leqslant
4 \reg(\ensddfv)^2 \poids{\diamant{D}}
\end{align*}
ce qui, avec \ref{eq:inter}, fournit le résultat.
\end{demo}

\noindent Les \Cref{prop:minoration_energie,prop:majoration_energie} aboutissent à l'inégalité
\[
\boxed{
\forall u^\ensddfv \in \Rddfvmn, 
\quad
\frac{1}{\sqrt{2} \reg(\ensddfv)^2} \, \frac{h_{\graphe^\ensddfv}}{\size(\ensddfv)} \normeddfvT{u^\ensddfv} \leqslant \normeddfvD{\grad^\ensdiamant u^\ensddfv}
}.
\]
Pour l'instant, la régularité $\reg(\ensddfv)$ contrôle la géométrie locale des mailles. Elle autorise cependant les raffinements locaux. Elle ne contrôle donc pas le rapport entre la plus grande et la plus petite taille de maille.

Dans la suite, nous allons démontrer que sous des hypothèses géométriques appropriées, il existe une constante $c > 0$ indépendante de $\size(\ensddfv)$ telle que
\[
 h_{\graphe^\ensddfv} \geqslant c \size(\ensddfv).
\]
Pour cela, nous allons ajouter une condition de quasi-uniformité globale du type
\[
\vartheta \size(\ensddfv) \leqslant \diam{\diamant{D}} \quad \forall \diamant{D} \in \ensdiamant.
\]

\begin{prop}
On suppose qu'il existe $N, c_A, C_\sigma > 0$ indépendantes de $\size(\ensddfv)$ telles que
\[
\degr(x_{\primal{K}}) \leqslant N,
\quad
\mesure(\primal{K}) \geqslant c_A \, \size(\ensddfv)^2,
\quad
\mesure(\sigma) \leqslant C_\sigma \, \size(\ensddfv).
\] 
Alors, pour tout $W \subset \vertices(\graphe^\ensprimal)$, on a
\[
\min\{ \vol(W), \vol(W^c) \} \leqslant \frac{N}{c_A \size(\ensddfv)^2} \min\{ \mesure(A_W), \mesure(\surflisse \setminus A_W) \},
\]
où
\[
A_W = \bigcup_{x_{\primal{K}} \in W} \primal{K} \subset \surflisse.
\]
\end{prop}
\begin{demo}
Puisque les mailles forment une partition,
\[
\mesure(A_W) = \sum_{x_{\primal{K}} \in W} \mesure(\primal{K}).
\]
Le volume combinatoire vaut
\[
\vol(W) = \sum_{x_{\primal{K}} \in W} \degr(x_{\primal{K}}).
\]
Or, pour chaque maille, 
\[
\degr(x_{\primal{K}}) \leqslant N \leqslant \frac{N}{c_A \size(\ensddfv)^2} \, \mesure(\primal{K}).
\]
En sommant, 
\[
\vol(W) \leqslant \frac{N}{c_A \size(\ensddfv)^2} \, \mesure(A_W).
\]
De même,
\[
\vol(W^c) \leqslant \frac{N}{c_A \size(\ensddfv)^2} \, \mesure(\surflisse \setminus A_W).
\]
On en déduit
\[
\min\{ \vol(W), \vol(W^c) \} \leqslant \frac{N}{c_A \size(\ensddfv)^2} \min\{ \mesure(A_W), \mesure(\surflisse \setminus A_W) \}.
\]
Pour obtenir $\mesure(\primal{K}) \geqslant c_A \size(\ensddfv)^2$, on peut combiner
\[
\frac{\diam{\primal{K}}}{\sqrt{\mesure(\primal{K})}} \leqslant \reg(\ensddfv)
\]
avec une quasi-uniformité globale
\[
\diam{\primal{K}} \geqslant \vartheta \size(\ensddfv).
\]
\begin{comment}
Une arête du bord de $A_S$ est une interface $\sigma = \arete{\primalplat{K}}{\primalplat{L}}$ telle que $x_{\primalplat{K}} \in S$ et $x_{\primalplat{L}} \not\in S$. 
\[
\partial A_S = \bigcup_{\segment{x_{\primalplat{K}}}{x_{\primalplat{L}}} \in \partial S} \sigma_{\primalplat{K} \primalplat{L}}.
\]
\end{comment}
\end{demo}

\begin{lemme}
\[
\abs{\partial W} \geqslant \frac{c}{\size(\ensddfv)} \mesure(\partial A_W).
\]
\end{lemme}
\begin{comment}
\begin{demo}
\[
\diam{\primal{K}} \geqslant \vartheta \size(\ensddfv).
\]
Une arête du bord de $A_S$ est une interface $\sigma = \arete{\primalplat{K}}{\primalplat{L}}$ telle que $x_{\primalplat{K}} \in S$ et $x_{\primalplat{L}} \not\in S$. 
\[
\partial A_S = \bigcup_{\segment{x_{\primalplat{K}}}{x_{\primalplat{L}}} \in \partial S} \sigma_{\primalplat{K} \primalplat{L}}.
\]
\end{demo}
\end{comment}

\begin{envide}[Stratégie]
Il faut alors établir deux comparaisons uniformes:
\[
\min\{ \vol(W), \vol(W^c) \} \leqslant \frac{C}{\size(\ensddfv)^2} \min\{ \mesure(A_W), \mesure(\surflisse \setminus A_W) \},
\]
et
\[
\abs{\partial W} \geqslant \frac{c}{\size(\ensddfv)} \mesure(\partial A_W).
\]
Elles donnent
\[
\frac{\abs{\partial W}}{\min\{ \vol(W), \vol(W^c) \}} \geqslant c \size(\ensddfv) \frac{\mesure(\partial A_W)}{\min\{ \mesure(A_W), \mesure(\surflisse \setminus A_W) \}} \geqslant c \size(\ensddfv) \, h(\surflisse).
\]
En prenant l'inf sur $W$, on obtient
\[
\boxed{
h_{\graphe^\ensprimal} \geqslant c \size(\ensddfv) \, h(\surflisse).
}
\]
Comme $\surflisse$ est une surface compacte connexe fixée, $h(\surflisse) > 0$, donc
\[
h_{\graphe^\ensprimal} \geqslant c_0 \size(\ensddfv).
\]
\end{envide}

\newpage

\begin{prop}
\[
\norme{\Pi^\variete{M}_\ensddfv u^\ensddfv}^2_{\L^2(\variete{M})} \sim h^2 \sum_{\primal{K} \in \ensprimal} u_{\primal{K}}{}^2
\]
\end{prop} 

\subsection{Passage du discret au continu}

\begin{defi}{Constante isopérimétrique}

\end{defi}

\begin{lemme}
L'inégalité \ref{eq:ineg_de_graphe} est vraie pour $h = h(\graphe)$.
\end{lemme}

\begin{lemme}
Soit $A \subset \variete{M}$,
\[
\inf \frac{\abs{\partial A}}{\abs{A}} \cdots
\]
\end{lemme}

\begin{lemme}
\[
    \frac{\abs{\partial S}}{\abs{S}} \sim \frac{\abs{\partial A}}{\abs{A}} \, h
\]
\end{lemme}

\begin{prop}[Relaxation condition de Cheeger]
    Pour $0 < k < 1$, on définit
    \[
    h_k(\variete{M}) \defeq \inf_{\substack{A \subset \variete{M} \\ \abs{A} \leqslant k \, \abs{\variete{M}}}} \frac{\abs{\partial A}}{\abs{A}},
    \]
    de sorte que $h(\variete{M}) = h_{1/2}(\variete{M})$. Alors
    \[
    \begin{cases}
        h_k(\variete{M}) = h(\variete{M}) &\text{si } \frac{1}{2} \leqslant k < 1 \text{ (relaxation)}, \\
        h_k(\variete{M}) \geqslant h(\variete{M}) &\text{si } 0 < k < \frac{1}{2} \text{ (restriction)}.
    \end{cases}
    \]
\end{prop}
\begin{demo}
On remarque que pour $0 < k \leqslant k' < 1$, alors $h_{k'}(\variete{M}) \leqslant h_k(\variete{M})$ car on minimise sur un ensemble plus grand. 

$h(\variete{M}) = h_{1/2}(\variete{M})$
\end{demo}

\begin{figure}[H]
    \centering
    \tikzset{
    ddfv/.style={
        scale=1.5,
        every path/.style={thick,line join=round, cap=round},
        ptsq/.style 2 args={
          rectangle,
          fill=##1,
          inner sep=2pt,
          label={[text=##1]##2}
        },
        ptsqbord/.style 2 args={
          rectangle,
          draw=##1,
          inner sep=2pt,
          label={[text=##1]##2}
        },
        ptcl/.style 2 args={
          circle,
          fill=##1,
          inner sep=1.5pt,
          label={[text=##1]##2}
        },
        ptclbord/.style 2 args={
          circle,
          draw=##1,
          inner sep=1.5pt,
          label={[text=##1]##2}
        },
        mailleprimale/.style={
            myblue,
            fill=mylightblue
        },
        mailledualeint/.style={
            myred,
            fill=mylightred
        },
        mailledualebord/.style={
            myred,
            fill=mylightred!50!white
        },
        maillediamantint/.style={
            mygreen,
            fill=mylightgreen
        },
        maillediamantbord/.style={
            mygreen,
            fill=mylightgreen!50!white
        }
    }
}

\tdplotsetmaincoords{60}{0}

%--------------------------------------------------
% Paramètres
%--------------------------------------------------
\def\R{1}
\def\Ntheta{6}
\def\Nphi{16}

% Sommets sélectionnés : couples i/j
\def\SelectedPoints{{2/12},{4/14},{5/11},{4/12}}

%--------------------------------------------------
% Tests de sélection
%--------------------------------------------------
\newif\ifselectedpoint
\newif\ifselectedtriangle

\newcommand{\testselectedpoint}[2]{%
  \selectedpointfalse
  \foreach \si/\sj in \SelectedPoints {%
    \ifnum#1=\si
      \ifnum#2=\sj
        \global\selectedpointtrue
      \fi
    \fi
  }%
}

\newcommand{\testselectedtriangle}[6]{%
  \selectedtrianglefalse
  \testselectedpoint{#1}{#2}%
  \ifselectedpoint
    \global\selectedtriangletrue
  \else
    \testselectedpoint{#3}{#4}%
    \ifselectedpoint
      \global\selectedtriangletrue
    \else
      \testselectedpoint{#5}{#6}%
      \ifselectedpoint
        \global\selectedtriangletrue
      \fi
    \fi
  \fi
}

%--------------------------------------------------
% Dessin d'un triangle plan
%--------------------------------------------------
\newcommand{\drawplanartriangle}[9]{%
  \testselectedtriangle{#1}{#2}{#3}{#4}{#5}{#6}%
  \ifselectedtriangle
    \draw[mailleprimale] #7 -- #8 -- #9 -- cycle;
  \else
    \draw[gray] #7 -- #8 -- #9 -- cycle;
  \fi
}

\newcommand{\drawplanartrianglegraphe}[9]{%
  \testselectedtriangle{#1}{#2}{#3}{#4}{#5}{#6}%
  \ifselectedtriangle
    \draw[mailleprimale, dashed] #7 -- #8 -- #9 -- cycle;
  \else
    \draw[gray, dashed] #7 -- #8 -- #9 -- cycle;
  \fi
}

\newcommand{\cheegerfig}[2][]{%
\begin{tikzpicture}[
  ddfv,
  scale=1.5,
  baseline=(current bounding box.center),
  #1
]
  \input{#2}
\end{tikzpicture}%
}

\begin{tikzpicture}[
  node distance=1.2cm,
  every node/.style={inner sep=0pt}
]

\node (G) {\cheegerfig{figures/cheeger_discret_to_continu_graphe.tex}};

\node[right=of G] (M) {\cheegerfig[tdplot_main_coords]{figures/cheeger_discret_to_continu_maillage.tex}};

\node[right=of M] (V) {\cheegerfig[tdplot_main_coords]{figures/cheeger_discret_to_continu_variete.tex}};

\draw[->, ultra thick, shorten >=4pt, shorten <=4pt] (G.east) -- (M.west);
\draw[->, ultra thick, shorten >=4pt, shorten <=4pt] (M.east) -- (V.west);

\end{tikzpicture}
    \caption{maillage non régulier}
\end{figure}



\begin{envide}[Sketch of the proof]
{\color{blue}
\begin{itemize}
    \item La courbure de la surface est minorée et donc dès que $h$ est suffisament petit, tous les triangles sont dans un voisinage tubulaire de la surface qui est tel qu'on peut toujours projeter de manière unique dessus (de tout point du maillage, on peut projeter). On peut projeter les maillages sur la surface et on peut le faire de sorte que les volumes restent en $h^2$ des triangles et les longueurs des arêtes restent en $h$ et peut être que les angles soient minorés par une constante uniforme indépendamment de $h$ (on peut l'avoir avec la contrainte sur la régularité du maillage).
    \item Les diamants doivent aussi être proches de la surface (il faudra projeter les diamants et estimer)
    \item Le $u$ est de deux sortes: $P^0$ du maillage duale. 
    \item On veut une inégalité comme:
    \[
    \norme{\widetilde{\grad} u}_{\L^2} \geqslant C \norme{u}_{\L^2}
    \]
    \item Tout le but de l'opération est de montrer que
    \begin{equation}\label{eq:inegdegraphe}
    \boxed{
    \sum_{i \sim j} (u_i - u_j)^2 \geqslant C \, h^2 \sum (u_i - c)^2
    }
    \end{equation}
    Il faut minimiser sur $c$, prendre la meilleure constante, la projection de $u$ sur les constantes.
    \item \ref{eq:inegdegraphe} est une inégalité de graphe, c'est le laplacien sur les graphe qui relie le membre de gauche au membre de droite avec une constante qui est la plus petite valeur propre du laplacien sur le graphe. \textbf{Il faut estimer cette constante}. 
    \item Ceci est vrai avec $h = h(G)$ (Cheeger) avec
    \[
    h(G) \defeq \inf_{\abs{S} \leqslant \frac{\abs{T_h}}{2}} \frac{\abs{\partial S}}{\abs{S}}
    \]
    \item Il y a autant d'arêtes qui relient $S$ aux autres que d'arêtes de bord qui contiennent les points de $S$.
    \item Tout les boulot est de montrer que 
    \[
    h(G) \geqslant h \, c
    \] 
    \item Soit $S$. On considère l'ensemble de tous les triangles qui sont de centre un point de $S$ et on projette ces triangles sur la surfaces. 
    \item On considère le sous-ensemble de la surface qui est le projeté des triangles du maillage pour lesquels il y a un point de $S$ qui est dedans. On appelle $A \subset \variete{M}$ cet ensemble. 
    \[
    A = \bigcup_{i \in S} \Delta^\variete{M}_i.
    \]
    \item D'après les lemmes de projection et les hypothèses : $\abs{\partial S} \sim \frac{\abs{\partial A}}{h}$ et $\abs{S} \sim \frac{\abs{A}}{h^2}$ et donc
    \[
    \frac{\abs{\partial S}}{\abs{S}} \sim \frac{\abs{\partial A}}{\abs{A}} \, h
    \]
    \item $\frac{\abs{\partial A}}{\abs{A}}$ : on est sur la surface, c'est la constante d'isopérimétrie de la variété et c'est minoré par quelque chose de strictement positif avec la condition $\boxed{\abs{A} \leqslant \frac{\abs{\variete{M}}}{2}}$.
\end{itemize}
}
\end{envide}

\subsection{État de l'art}